\chapter{Deciding Under Uncertainty}\label{ch:decisions-under-uncertainty}

% Scope: expected value, Bayes, decision trees, cognitive bias.
% Cross-link: Bayesian updating is reused by data-driven attribution & MMM (\cref{ch:attribution-and-measurement});
%   the chain rule of probability underlies the conversion funnel (\cref{ch:digital-funnel-and-crm}).

Every business decision is a bet placed before the outcome is known. The frameworks
here --- expected value, Bayes, and prospect theory --- do not remove uncertainty; they make the
structure of the bet explicit so the manager can evaluate it clearly rather than react to it
emotionally.

\section{Expected Value and Decision Trees}\label{sec:expected-value}
% complexity: 4 -> figure decision-tree

When two options have different risk profiles, which do you choose? Expected value converts
probability-weighted outcomes into a single comparable number, but a decision tree forces you
to make the structure of the uncertainty explicit before you compute. Comparing a safe outcome
with a risky one by gut feel systematically misleads because humans weight recent, vivid, and
emotionally salient outcomes too heavily. Expected value is the long-run arithmetic average
payoff --- the number you would converge to if the decision were made many times under identical
conditions. For one-off decisions, EV is a disciplined baseline, not a guarantee.

For a discrete set of mutually exclusive outcomes \(x_i\) with probabilities \(p_i\)
(\(\sum_i p_i = 1\)):
\begin{equation}\label{eq:ev}
  \mathbb{E}[X] \;=\; \sum_i p_i\, x_i.
\end{equation}
A decision tree maps each choice to a set of probability-weighted branches, with EV computed at
each terminal node and folded back to the decision node by \cref{eq:ev}. The tree's value is
that it forces the analyst to enumerate outcomes and assign probabilities rather than reason
from a single expected scenario.

\begin{figure}[htbp]
  \centering
  \includegraphics[width=0.86\linewidth]{decision-tree}
  \caption{Decision tree for two Target CPA bids. The higher bid (\money{900}) has higher expected
    value despite a lower success probability, because the downside is smaller relative to the
    upside. Rolling back the tree to the decision node reveals this only after probabilities are
    made explicit.}
  \label{fig:decision-tree}
\end{figure}

\Cref{eq:ev} assumes outcomes and probabilities are known. In practice, both require estimation
from data or expert judgment, introducing model risk. EV also ignores the distribution around the
mean: two options with identical EV can have very different variance, and a firm with constrained
capital may rationally prefer the lower-variance option even at a lower EV (risk aversion ---
formalised in \cref{sec:prospect-theory}). Finally, EV from a decision tree collapses to a single
number and discards the timing of cash flows; embed discounting when comparing multi-period paths.

\begin{example}
The running SaaS business considers two Target CPA levels for a new campaign. At
\(\mathrm{CPA}=\money{600}\), the conversion model assigns a \qty{70}{\percent} probability of
acquiring a customer worth \(\mathrm{LTV}=\money{3150}\) and \qty{30}{\percent} of a
low-engagement customer worth \money{900} (three months then churned):
\[
  \begin{aligned}
    \mathbb{E}[\text{value} \mid \mathrm{CPA}=\money{600}]
      &= 0.70\times(\money{3150}-\money{600}) + 0.30\times(\money{900}-\money{600}) \\
      &= \money{1785} + \money{90} = \money{1875}.
  \end{aligned}
\]
At \(\mathrm{CPA}=\money{900}\), the higher spend attracts a higher-intent segment:
\qty{85}{\percent} chance of the full LTV customer, \qty{15}{\percent} of the low-engagement one:
\[
  \begin{aligned}
    \mathbb{E}[\text{value} \mid \mathrm{CPA}=\money{900}]
      &= 0.85\times(\money{3150}-\money{900}) + 0.15\times(\money{900}-\money{900}) \\
      &= \money{1912.50}.
  \end{aligned}
\]
The \money{900} CPA has higher expected per-customer profit (\money{1912.50} vs.\ \money{1875})
despite spending \money{300} more per acquisition. The decision: raise the Target CPA. The trap
is choosing \money{600} because it \emph{feels} more conservative --- \cref{sec:prospect-theory}
explains why that intuition is systematically wrong.
\end{example}

\begin{admonition}{Warning}
Single-scenario planning --- building one set of projections, labelling it ``base case,'' and
evaluating it as if it were certain --- is a decision tree with one branch, which assigns
probability~1 to the chosen scenario. Making the branches and their probabilities explicit is not
additional work; it is the minimum discipline that EV analysis requires \parencite{kahneman2011}.
\end{admonition}

When decisions are sequential and information arrives between stages, the tree becomes a real
options problem: the value of being able to wait, expand, or abandon is an option premium on top
of the static EV. Black-Scholes pricing is the continuous-time limit of a binomial tree. The
practical implication: staging capital commitments (first a pilot, then rollout) is almost always
worth more than its cost, because it preserves the option to abandon. Modelling this requires
option-pricing methods, not plain EV (\cref{ch:valuation}).

Expected value is the language of budgeting and scenario analysis; it feeds directly into the
DCF framework of \cref{ch:valuation} and the MMM attribution models of
\cref{ch:attribution-and-measurement}. \textcite{anderson2019statistics} cover the probability
machinery; \textcite{kahneman2011} document why managers deviate from it.

\section{Bayesian Updating}\label{sec:bayesian-updating}
% complexity: 3

A new signal has arrived --- a whitepaper download, a pilot result, a churn signal. How much
should it move your prior belief? Bayes's theorem gives the exact answer; the practical skill is
choosing what counts as evidence and calibrating its diagnostic power. Frequentist inference asks
what the probability of this data is given a fixed true parameter; Bayesian inference asks what
the probability the parameter is true is, given the data just observed. The second question is
almost always what the manager needs. Bayes says: start with a prior (your initial belief),
observe evidence, and update proportionally to how much more likely the evidence is under the
hypothesis than under its alternatives.

Let \(H\) be a hypothesis and \(D\) observed data. The posterior probability of \(H\) given \(D\) is:
\begin{equation}\label{eq:bayes}
  P(H \mid D) \;=\; \frac{P(D \mid H)\,P(H)}{P(D)},
\end{equation}
where \(P(H)\) is the prior, \(P(D \mid H)\) the likelihood of the data under \(H\), and
\(P(D)=P(D\mid H)P(H)+P(D\mid \neg H)P(\neg H)\) the normalising constant (marginal likelihood).
The update is multiplicative in log-odds: \(\log\frac{P(H\mid D)}{P(\neg H\mid D)} =
\log\frac{P(H)}{P(\neg H)} + \log\frac{P(D\mid H)}{P(D\mid \neg H)}\). Each new observation
adds its log-likelihood ratio to the running log-odds.

\Cref{eq:bayes} is exact but only as good as its inputs. A miscalibrated prior (e.g.\ a base rate
ignored in favour of the vivid case) propagates into the posterior. A likelihood built on
non-representative historical data likewise distorts the update. The formula cannot fix bad inputs;
it can only make the structure of the reasoning transparent so that the inputs can be challenged.

\begin{example}
The sales team believes \qty{25}{\percent} of trial users eventually convert to paid --- the prior:
\(P(\text{paid})\!=\!0.25\). A whitepaper download is observed. Historical data show that
\qty{60}{\percent} of eventually-paid users downloaded a whitepaper during trial, versus only
\qty{15}{\percent} of users who churned: \(P(D\mid\text{paid})=0.60\),
\(P(D\mid\text{not paid})=0.15\). Applying \cref{eq:bayes}:
\[
  P(D) = 0.60\times0.25 + 0.15\times0.75 = 0.15 + 0.1125 = 0.2625,
\]
\[
  P(\text{paid}\mid D) = \frac{0.60\times0.25}{0.2625} \approx 0.571.
\]
A single whitepaper download more than doubles the estimated conversion probability from \qty{25}{\percent}
to \qty{57}{\percent}. The CRM implication: route this lead to a higher-touch sales sequence whose
cost is justified by the revised LTV estimate. Without Bayesian updating, both leads receive
identical treatment despite very different expected values.
\end{example}

\begin{admonition}{Warning}
Base-rate neglect is the most common Bayesian error: updating on the vivid signal while ignoring
the prior. A startup announces its trial-to-paid rate is \qty{40}{\percent} --- but is that
conditional on having the same acquisition funnel and product quality as your benchmark? Without
a calibrated prior and a carefully measured likelihood ratio, the \qty{40}{\percent} benchmark
contaminates rather than informs the posterior \parencite{kahneman2011}.
\end{admonition}

Bayesian A/B testing replaces the frequentist null-hypothesis framework with a posterior over the
conversion-rate difference. This yields a decision-relevant quantity (``what is the probability
that variant B is better?'') rather than a p-value (``how likely are these data if the null is
true?''). Multi-armed bandit algorithms extend the framework to continuously reallocate traffic
toward better-performing variants, making the explore--exploit trade-off explicit. The
attribution models in \cref{ch:attribution-and-measurement} are Bayesian at their core:
they update channel-attribution weights as new conversion data arrive.

Bayesian updating is the backbone of probabilistic attribution in \cref{ch:attribution-and-measurement}
and of the customer-scoring models that feed the funnel in \cref{ch:digital-funnel-and-crm}.
\textcite{anderson2019statistics} cover the frequentist and Bayesian frameworks in parallel.

\section{Cognitive Bias and Prospect Theory}\label{sec:prospect-theory}
% complexity: 4 -> figure prospect-value-fn

Your team has two options with identical expected value. They will reliably choose the less risky
one. Prospect theory explains the specific, predictable ways that human risk perception deviates
from EV maximisation, so you can correct for it. People evaluate outcomes relative to a reference
point (usually the status quo), not in absolute terms. Gains feel good but with diminishing
sensitivity; losses from the same reference point feel roughly twice as bad as equivalent gains
feel good. The resulting value function is S-shaped: concave in the gain domain (risk-averse),
steeply convex in the loss domain (risk-seeking to avoid a certain loss). This is why a certain
\money{600} CAC ``feels safer'' than a \money{900} bid with higher expected profit: the
\money{300} extra spend is coded as a certain loss, weighted at roughly double its face value.

\textcite{kahneman1979prospect} specify the prospect theory value function as:
\begin{equation}\label{eq:prospect}
  v(x) \;=\;
  \begin{cases}
    x^{\alpha} & x \ge 0, \\[3pt]
    -\lambda\,(-x)^{\beta} & x < 0,
  \end{cases}
\end{equation}
where \(x\) is the outcome relative to the reference point, \(\alpha,\beta\in(0,1)\) capture
diminishing sensitivity in both domains (\(\alpha\approx\beta\approx0.88\) empirically), and
\(\lambda>1\) is the \emph{loss-aversion coefficient} (\(\lambda\approx2.25\) on average). The
function is not a utility function over wealth; it is a value function over changes.
Prospect theory also replaces objective probabilities with decision weights that overweight
small probabilities and underweight moderate-to-large ones, but the value function drives most
of the managerially relevant predictions.

\begin{figure}[htbp]
  \centering
  \includegraphics[width=0.86\linewidth]{prospect-value-fn}
  \caption{The prospect theory value function (\cref{eq:prospect}, \(\lambda=2.25\)). Concave in gains
    (risk-averse); steeply convex in losses (risk-seeking to avoid certain loss). The asymmetry
    at zero means a loss of any given size looms larger than an equal gain.}
  \label{fig:prospect-value-fn}
\end{figure}

\Cref{eq:prospect} is a descriptive model of average human behaviour under laboratory conditions,
not a normative rule. It does not hold uniformly across cultures, stakes levels, or domains with
repeated-play experience. The reference point itself shifts: a bonus framed as avoiding a
penalty produces different choices than the same bonus framed as a gain, even when the objective
outcome is identical (the \emph{framing effect}). Managers who understand this can design
compensation structures and experiment reports that reduce bias, but they should not assume
\(\lambda=2.25\) is a universal constant.

\begin{example}
Return to the two CPA bids. The \money{600} CPA feels safer because paying \money{900} is coded
as a \money{300} additional loss relative to the \money{600} reference point. Applying
\cref{eq:prospect} with \(\lambda=2.25\):
\[
  v(-\money{300}) \approx -2.25\times(300)^{0.88} \approx -2.25\times191 \approx -430
  \quad\text{(loss)}
\]
versus the gain from higher expected return:
\[
  v(+\money{37.50}) \approx (37.50)^{0.88} \approx 27
  \quad\text{(gain from EV difference \money{1912.50} vs.\ \money{1875})}.
\]
The psychic cost of the larger bid (\(-430\) value units) dwarfs the psychic gain of the higher
EV (+27 units) --- so the manager rationally feels the \money{600} bid is better, even though
expected profit is higher at \money{900}. Correct for it by anchoring the decision to EV
(\cref{eq:ev}), not to the magnitude of the incremental spend.
\end{example}

\begin{admonition}{Warning}
The planning fallacy --- underestimating the cost and duration of a project because the reference
point is the plan, not historical base rates for similar projects --- compounds loss aversion:
managers are simultaneously overconfident about outcomes and excessively averse to downside
scenarios they cannot vividly imagine \parencite{tversky1974judgment}. Pre-mortems, in which the
team imagines the project has already failed and asks why, are the most reliably effective
debiasing technique.
\end{admonition}

Prospect theory's normative status is contested. \textcite{kahneman1979prospect} present it as a
descriptive model; expected utility theorists (\cref{eq:ev}) argue that deviation from EV
maximisation is simply a decision error that better incentives or education can eliminate.
Behavioural economists respond that the biases are too stable and too large to be mere noise.
The practical resolution: use EV as the decision criterion, but design the choice architecture ---
how options are framed, sequenced, and presented --- to reduce the gap between EV and
prospect-theory evaluations \parencite{kahneman2011}.

\textcite{thaler1985mental} introduces mental accounting --- the tendency to evaluate outcomes in
separate psychological accounts rather than as changes in total wealth. A marketing implication:
bundling a price increase with a visible product improvement moves the net change from the loss
account to the ``fair exchange'' account, reducing the perceived loss. Framing effects and mental
accounting are the mechanism behind anchoring (\cref{sec:prospect-theory} worked example) and
price partitioning strategies explored in \cref{ch:pricing}.

Prospect theory is the foundation for pricing architecture in \cref{ch:pricing} and for
understanding why customers resist switching even when the alternative has higher expected value.
The Bayesian and prospect-theoretic frameworks jointly underpin the attribution discussion in
\cref{ch:attribution-and-measurement}. \textcite{tversky1974judgment} document the heuristics;
\textcite{kahneman1979prospect} formalise the value function; \textcite{kahneman2011} synthesise
both into a practitioner framework.
